%% parts/p04/degrees-and-priority.tex
%% Migrates v1.1 lines 8898--9079 and fills the registered union
%% D-P5-003/004/005 + N-a-016/017/018.  Source statuses are unchanged.

\chapter[Relative Computation and Degrees]{Relative Computation, Degrees, and Priority}
\label{sec:p04-degrees-priority}

Many-one reducibility translates one membership question into another before
the target computation begins.  Relative computation instead permits a
machine to consult a fixed set during its run.  The distinction supports both
the jump hierarchy and the finer partial order of computational degrees.

\section{Oracle computation and the jump}
\label{sec:p04-oracles-jump}

\begin{hbdefinition}{def:p04-oracle-reducibility}{Oracle computation and Turing reducibility}
An \emph{oracle Turing machine with oracle \(B\subseteq\N\)} is a
deterministic single-tape machine with a query operation: after computing a
number \(n\), it enters one of two designated states according as \(n\in B\)
or \(n\notin B\).  Its numerical input, acceptance, rejection, valid-output,
and divergence conventions are exactly those for ordinary machines.

A set \(A\subseteq\N\) is \emph{Turing reducible} to \(B\), written
\(A\leq_TB\), when an oracle machine with oracle \(B\) decides \(A\).  A set
is \emph{computably enumerable relative to \(B\)} when it is the domain of a
partial numerical function computed by an oracle machine with oracle \(B\).
\index{oracle Turing machine}
\index{Turing reducibility}
\index{computably enumerable!relative}
\end{hbdefinition}

\begin{hblemma}{thm:p04-many-one-versus-turing}{proved=annex:p4.degrees.5}{Many-one versus Turing reducibility}
Every many-one reduction gives a one-query Turing reduction.  The converse
fails; specifically,
\[
  \N\setminus K\leq_TK
  \qquad\text{but}\qquad
  \N\setminus K\not\leq_mK.
\]
\end{hblemma}

Fix the source numbering of oracle machines: it uses the ordinary
normalization of designated slots, includes query rows, and makes invalid
codes diverge on every input.  Write \(\varphi^B_e(x)\) for the numerical output of
the valid accepting computation, leaving it undefined on rejection,
divergence, or invalid accepting output.

\begin{hbdefinition}{def:p04-turing-jump}{Turing jump}
The \emph{Turing jump} of \(B\subseteq\N\) is
\[
 B'\defeq\left\{e\in\N:\varphi^B_e(e)\mathord\downarrow\right\}.
\]
Set \(B^{(0)}\defeq B\) and \(B^{(n+1)}\defeq(B^{(n)})'\); in particular,
\(\varnothing^{(n)}\) is the \(n\)-th jump of the empty set.  Definedness,
not arbitrary halting, is the convergence notion fixed by
Convention~\ref{def:p04-convergence-notation}; this is the R-004 repair.
\index{Turing jump}
\end{hbdefinition}

\begin{hblemma}{thm:p04-relativized-adequacy}{sketched=annex:p4.degrees.1}{Relativized adequacy}
For every oracle \(B\subseteq\N\):
\begin{enumerate}
 \item every partial function computed by an oracle machine with oracle \(B\)
       is \(\varphi^B_e\) for some valid index \(e\);
 \item one fixed oracle machine computes
       \(\mathsf{Univ}^B(e,x)\simeq\varphi^B_e(x)\); and
 \item for all \(m,n\geq1\), a total computable function
       \(\mathsf{par}_n^m\), independent of \(B\), satisfies
       \[
        \varphi^B_{\mathsf{par}_n^m(e,\vec a)}(\vec x)
        \simeq\varphi^B_e(\vec a,\vec x)
       \]
       simultaneously for every oracle and all displayed arguments.
\end{enumerate}
\end{hblemma}

Consequently, relative c.e. sets are exactly the domains of the functions
\(\varphi^B_e\).

\begin{hbtheorem}{thm:p04-jump}{sketched=annex:p4.degrees.2}{Jump theorem}
For every \(B\subseteq\N\), the jump \(B'\) is computably enumerable
relative to \(B\), but \(B'\not\leq_TB\).
\end{hbtheorem}

\begin{hbtheorem}{thm:p04-post}{cited=Shoenfield1967}{Post's theorem}
For every \(n\in\N\) and \(A\subseteq\N\):
\begin{enumerate}
 \item \(A\in\boldsymbol\Sigma^0_{n+1}\) if and only if \(A\) is
       computably enumerable relative to \(\varnothing^{(n)}\); and
 \item \(A\in\boldsymbol\Delta^0_{n+1}\) if and only if
       \(A\leq_T\varnothing^{(n)}\).
\end{enumerate}
\end{hbtheorem}

The prefix-form transfer theorem licenses the formulation in the handbook's
closure classes.  See also \cite{Barwise1977}.  Together with the jump
theorem, Post's theorem gives the separation of consecutive levels; the
remaining join/completeness steps belong to the cited proof.

\begin{hblemma}{thm:p04-jump-convention-bridge}{sketched=annex:p4.degrees.3}{Convention bridge for the jump}
One has \(K\sim_m\varnothing'\): the source's definedness-based jump is
many-one equivalent to the usual halting-based jump, and jump monotonicity
holds:
\[
 A\leq_TB\quad\text{implies}\quad A'\leq_mB'.
\]
Moreover \(B\leq_mB'\) and \(B'\not\leq_TB\), so \(B\) is strictly
Turing reducible to \(B'\).  Thus the cited form of Post's theorem is invariant
under these two jump conventions.
\end{hblemma}

\begin{table}[htbp]
\centering
\scriptsize
\begin{tabular}{@{}p{.12\linewidth}p{.27\linewidth}p{.24\linewidth}p{.27\linewidth}@{}}
\toprule
Class & Normal form and closure & Computability reading & Calibrating set/result \\
\midrule
\(\Sigma^0_0=\Pi^0_0\) & bounded negation-normal-form matrix; closed under
Boolean operations and marked bounded quantifiers & bounded-formula level & no
completeness claim made here \\
\(\boldsymbol\Sigma^0_1\) & existential prefix; closed under
\(\land,\lor\), marked bounded quantifiers, and existential quantification &
c.e. sets & \(K\) is many-one complete \\
\(\boldsymbol\Pi^0_1\) & universal-prefix dual of the preceding row &
complements of c.e. sets & \(\N\setminus K\) is many-one complete \\
\(\boldsymbol\Delta^0_1\) & both level-one forms & decidable sets &
Theorem~\ref{thm:p04-ce-sigma-one} \\
\(\boldsymbol\Pi^0_2\) & \(\forall\vec u\exists\vec v\,\psi\), bounded
\(\psi\) & universal--existential level & \(\mathsf{TOT}\) is many-one
complete \\
\(\boldsymbol\Sigma^0_{n+1}\) & \(n+1\) alternating blocks, existential
first; the closure clauses are those of
Theorem~\ref{thm:p04-prefix-normal-form} & c.e. relative to
\(\varnothing^{(n)}\) & \(\varnothing^{(n+1)}\) is the standard complete
representative \cite{Shoenfield1967} \\
\(\boldsymbol\Pi^0_{n+1}\) & universal-prefix dual & complements of sets c.e.
relative to \(\varnothing^{(n)}\) & the complement of
\(\varnothing^{(n+1)}\) is complete \cite{Shoenfield1967} \\
\(\boldsymbol\Delta^0_{n+1}\) & both level-\(n+1\) forms & exactly the sets
Turing reducible to \(\varnothing^{(n)}\) &
Theorem~\ref{thm:p04-post}, for \(n\geq0\) \\
\bottomrule
\end{tabular}
\caption{Arithmetical hierarchy (apparatus A-047).  Formula classes and
boldface relation classes remain distinct; the prefix/closure equivalence is
Theorem~\ref{thm:p04-prefix-normal-form}, and every completeness claim is
under many-one reductions.}
\label{tab:p04-arithmetical-hierarchy}
\end{table}

\section{Turing degrees}
\label{sec:p04-turing-degrees}

\begin{hbdefinition}{def:p04-turing-degree}{Turing equivalence and degrees}
The preorder \(\leq_T\) compares sets but does not identify different sets
with the same computational power.  Define
\[
 A\sim_TB\quad\Leftrightarrow\quad
 A\leq_TB\ \text{ and }\ B\leq_TA.
\]
The \emph{Turing degree} of \(A\) is its equivalence class
\(\deg_T(A)=[A]_T\).  The set of degrees is ordered by
\[
 [A]_T\leq[B]_T\quad\Leftrightarrow\quad A\leq_TB.
\]
This is well defined and antisymmetric precisely because representatives are
quotiented by mutual reducibility.  A degree therefore names an equivalence
class of sets, not a set, a machine, or a reducibility statement.  The
distinction is used by join, jump, and Post's problem below.
\index{Turing degree}
\index{Turing equivalence}
\end{hbdefinition}

\begin{hbdefinition}{def:p04-degree-operations}{Join and jump on degrees}
For \(A,B\subseteq\N\), put
\[
 A\oplus B=\{2n:n\in A\}\cup\{2n+1:n\in B\}.
\]
Define the join of \(\deg_T(A)\) and \(\deg_T(B)\) to be
\(\deg_T(A\oplus B)\).  Lemma~\ref{thm:p04-degree-semilattice} verifies that
this operation is independent of representatives and is their least upper
bound.  The monotonicity in
Lemma~\ref{thm:p04-jump-convention-bridge} makes
\[
 \deg_T(A)'\defeq\deg_T(A')
\]
well defined.  Write \(\mathbf0=\deg_T(\varnothing)\) and
\(\mathbf0'=\deg_T(\varnothing')\).  A \emph{c.e. degree} is the degree of
some computably enumerable set.  These are derived operations and subclasses
of the already defined degree order, not additional primitive relations.
\index{Turing degree!join}
\index{Turing degree!jump}
\index{computably enumerable degree}
\end{hbdefinition}

\begin{hblemma}{thm:p04-degree-semilattice}{cited=KleenePost1954}{Upper-semilattice structure of the Turing degrees}
The order on Turing degrees and the join operation in
Definition~\ref{def:p04-degree-operations} are independent of representatives.
The degree \(\deg_T(A\oplus B)\) is an upper bound of
\(\deg_T(A)\) and \(\deg_T(B)\), and it is below every common upper bound.
\end{hblemma}

\begin{hbtheorem}{thm:p04-kleene-post}{cited=KleenePost1954}{Kleene--Post incomparability theorem}
There are Turing degrees \(\mathbf a,\mathbf b<\mathbf0'\) such that neither
\(\mathbf a\leq\mathbf b\) nor \(\mathbf b\leq\mathbf a\).  No claim that
these two degrees are c.e. is made by this theorem.
\end{hbtheorem}

\section{Post's problem and finite injury}
\label{sec:p04-finite-injury}

Post's problem asks whether some c.e. degree lies strictly between
\(\mathbf0\) and \(\mathbf0'\).  Friedberg and Muchnik answered it
affirmatively by constructing incomparable c.e. sets
\cite{Friedberg1957,Muchnik1956}.

\begin{hbmethod}{def:p04-finite-injury-method}{Finite-injury priority construction}
Construct c.e. sets \(A\) and \(B\) in effective stages while meeting, in the
order
\[
 R_0>S_0>R_1>S_1>\cdots,
\]
the requirements
\[
 R_e:\varphi_e^B\not\simeq\operatorname{char}_A,
 \qquad
 S_e:\varphi_e^A\not\simeq\operatorname{char}_B.
\]
Here \(\varphi_e^B\) is the partial function computed by the \(e\)-th
program in the fixed acceptable oracle-machine numbering, while
\(\operatorname{char}_A\) and \(\operatorname{char}_B\) are the total
characteristic functions of the constructed sets.  The requirements therefore
compare functions extensionally rather than comparing a function with a set.
An \(R_e\)-strategy reserves a fresh witness \(x\) above all active
higher-priority restraints and distinct from every reserved witness, keeps
\(x\notin A\), and waits for a convergent computation
\(\varphi_e^B(x)=0\).  Its query use \(u\) is the largest oracle position
queried by that finite computation; if it makes no query, its oracle restraint
is empty.  When the computation appears, enumerate \(x\) into \(A\), preserve
every oracle answer below \(u+1\), and declare the requirement satisfied.
Computations with value \(1\) require no enumeration action, and the strategy
continues waiting when no convergent computation appears.  The
\(S_e\)-strategy is symmetric.

If a higher-priority strategy later enumerates below a lower strategy's
oracle restraint, initialize the lower strategy again: cancel its witness and
restraint, but never remove a number already enumerated; then choose a new
fresh witness.  No lower-priority strategy may enumerate within a protected
oracle segment unless a higher-priority action first injures and initializes
it.  This is a construction scheme for enumerations of \(A,B\), not a new
reducibility or a new kind of degree.  Annex~\ref{annex:p4.degrees.4} records
the priority induction, satisfaction verification, and remaining bookkeeping
omissions.
\index{finite-injury priority method}
\index{priority argument}
\end{hbmethod}

\begin{hbtheorem}{thm:p04-friedberg-muchnik}{sketched=annex:p4.degrees.4}{Friedberg--Muchnik theorem}
There are computably enumerable sets \(A,B\subseteq\N\) such that
\(A\not\leq_TB\) and \(B\not\leq_TA\).  Consequently there are incomparable
c.e. degrees strictly between \(\mathbf0\) and \(\mathbf0'\), answering
Post's problem.
\end{hbtheorem}

\begin{table}[htbp]
\centering
\small
\begin{tabular}{@{}p{.65in}p{1.15in}p{1.25in}p{1.25in}p{1.30in}@{}}
\toprule
Strategy & Witness policy & Wait condition & Positive action & Restraint and injury \\
\midrule
\(R_e\) & reserve fresh \(x\notin A\) & \(\varphi_e^B(x)\downarrow=0\) & enumerate \(x\) into \(A\) & preserve every \(B\)-answer below \(u+1\); reset after higher-priority injury \\
\(S_e\) & reserve fresh \(y\notin B\) & \(\varphi_e^A(y)\downarrow=0\) & enumerate \(y\) into \(B\) & preserve every \(A\)-answer below \(u+1\); reset after higher-priority injury \\
divergence & retain the final witness & no convergence ever appears & none & permanent disagreement by undefinedness \\
value \(1\) & keep the witness out of the target set & no stable value \(0\) appears & none & after oracle stabilization, a final value \(1\) disagrees with nonmembership \\
verification & induct on priority & higher strategies have stabilized & final witness is chosen & only finitely many resets; final restraint survives \\
\bottomrule
\end{tabular}
\caption{Finite-injury construction table for Friedberg--Muchnik
(N-a-017/D-P5-004).  The table records the construction, while
Theorem~\ref{thm:p04-friedberg-muchnik} and its annex item carry the evidence
status.}
\label{tab:p04-finite-injury}
\end{table}

\section{Structure of the c.e. sets}
\label{sec:p04-ce-structure}

\begin{hbconvention}{def:p04-ce-numbering}{Acceptable numbering of the c.e. sets}
The index-sensitive results below require more than an arbitrary effective
enumeration.  For this chapter, an \emph{acceptable numbering of programs}
means a numbering satisfying the enumeration/surjectivity clause~(b) of
Theorem~\ref{thm:p04-enumeration}, the universal property of
Theorem~\ref{thm:p04-universal-computation}, and the uniform parameter
property of Theorem~\ref{thm:p04-parameter}.  The fixed program numbering has
exactly this package.  Put
\[
 W_e\defeq\{x\in\N:\varphi_e(x)\mathord\downarrow\}.
\]
Thus \((W_e)_{e\in\N}\) is the induced acceptable numbering of the c.e.
subsets of \(\N\).  The adjective \emph{acceptable} here names a
metalinguistic package of numbering hypotheses, not a reducibility, a degree,
or a new class of sets; unlike a merely effective enumeration, it supplies
the uniform index translations used in productivity, creativity, Myhill's
theorem, and Rice--Shapiro.
\index{acceptable numbering}
\index{computably enumerable set!acceptable numbering}
\end{hbconvention}

\begin{hbdefinition}{def:p04-one-one-reduction}{One-one reducibility}
A \emph{one-one reduction} of \(A\subseteq\N\) to \(B\subseteq\N\) is a
total computable injective function \(f\) such that
\[
 x\in A\quad\Leftrightarrow\quad f(x)\in B.
\]
Write \(A\leq_1B\).  This strengthens many-one reducibility by requiring the
translator to be injective; the distinction is used by Myhill's theorem.
\index{one-one reducibility}
\end{hbdefinition}

\begin{hbdefinition}{def:p04-simple-set}{Simple set}
A c.e. set \(S\) is \emph{simple} when \(\N\setminus S\) is infinite but
contains no infinite c.e. subset.  Simplicity is a property of a c.e. set,
not a completeness notion: its use is to witness a noncomputable c.e. set.
Simplicity alone neither implies nor excludes Turing completeness.
\index{simple set}
\end{hbdefinition}

\begin{hbtheorem}{thm:p04-simple-set-exists}{cited=Post1944}{Existence of simple sets}
There exists a simple computably enumerable set.  Every simple set is
noncomputable.
\end{hbtheorem}

\begin{hbdefinition}{def:p04-productive-creative}{Productive complements and creative sets}
A set \(C\) is \emph{productive} when some partial computable function \(p\)
has the following property: whenever \(W_e\subseteq C\), the value \(p(e)\)
is defined and lies in \(C\setminus W_e\).  A c.e. set \(A\) is
\emph{creative} when its complement is productive.
Productivity supplies a uniform escape from every c.e. subset; creativity
combines that property with enumerability and is used in the completeness
theorem below.
\index{productive set}
\index{creative set}
\end{hbdefinition}

\begin{hbtheorem}{thm:p04-creative-complete}{cited=Myhill1955}{Creative-set characterization}
A computably enumerable set is creative if and only if it is
many-one complete among the computably enumerable sets.  In particular, the
self-application domain set \(K\) is creative.
\end{hbtheorem}

\begin{hbtheorem}{thm:p04-myhill}{cited=Myhill1955}{Myhill isomorphism theorem}
If \(A\leq_1B\) and \(B\leq_1A\), then some computable permutation
\(\pi:\N\to\N\) satisfies
\[
 x\in A\quad\Leftrightarrow\quad\pi(x)\in B.
\]
Consequently, any two creative sets are computably isomorphic.
\end{hbtheorem}

\begin{hbtheorem}{thm:p04-rice-shapiro}{cited=Johnson1977Rice}{Rice--Shapiro theorem}
For \(d\in\N\), let \(D_d\) be the finite subset of \(\N\) whose members are
the positions of the \(1\)-bits in the binary expansion of \(d\).  Let
\(\mathcal C\) be a class of c.e. subsets of \(\N\), and let
\(I_{\mathcal C}=\{e:W_e\in\mathcal C\}\) be its extensional index set.
Then \(I_{\mathcal C}\) is computably enumerable if and only if there is a
computably enumerable set \(\mathsf{Fin}_{\mathcal C}\subseteq\N\) of
finite-set codes such
that, for every c.e. set \(W\),
\[
 W\in\mathcal C
 \quad\Leftrightarrow\quad
 \exists d\in\mathsf{Fin}_{\mathcal C}\;(D_d\subseteq W).
\]
The empty class is included by allowing
\(\mathsf{Fin}_{\mathcal C}=\varnothing\).
Thus a positive c.e. extensional property is witnessed by finite positive
information.  The result concerns extensional index sets of c.e. sets, not
arbitrary semantic properties of programs; within that scope it is stronger
than Rice's undecidability classification.
\end{hbtheorem}

\begin{table}[htbp]
\centering
\small
\begin{tabular}{@{}p{1.12in}p{1.25in}p{1.30in}p{1.55in}@{}}
\toprule
Object/fact & Exact standing & Order-theoretic content & Evidence \\
\midrule
Turing degrees & \(\mathcal P(\N)/{\sim_T}\) & partial order induced by \(\leq_T\) & Definition~\ref{def:p04-turing-degree} \\
join & \(\deg_T(A\oplus B)\) & least upper bound of two degrees & Lemma~\ref{thm:p04-degree-semilattice}; \cite{KleenePost1954} \\
jump & \(\deg_T(A)'=\deg_T(A')\) & strictly increases every degree & Theorem~\ref{thm:p04-jump} and Lemma~\ref{thm:p04-jump-convention-bridge} \\
Kleene--Post & arbitrary degrees below \(\mathbf0'\) & incomparable pair exists & \cite{KleenePost1954} \\
Post's problem & c.e. degrees between \(\mathbf0\) and \(\mathbf0'\) & asks for a nontrivial intermediate c.e. degree & \cite{Post1944} \\
Friedberg--Muchnik & c.e. degrees & incomparable pair, hence intermediate degrees & sketched annex; \cite{Friedberg1957,Muchnik1956} \\
\bottomrule
\end{tabular}
\caption{Turing degrees and jump facts (apparatus A-049).  The word
\emph{arbitrary} in the
Kleene--Post row prevents the false inference that its witnesses are the c.e.
witnesses later produced by finite injury.}
\label{tab:p04-degree-facts}
\end{table}

\begin{figure}[htbp]
\centering
\begin{tikzpicture}[
  node distance=14mm and 18mm,
  degree/.style={circle,draw,minimum size=9mm,font=\small},
  order/.style={-{Latex[length=2mm]},thick},
  jumpop/.style={-{Latex[length=2mm]},thick,dashed}
]
 \node[degree] (zero) {\(\mathbf0\)};
 \node[degree,above left=of zero] (a) {\(\mathbf a\)};
 \node[degree,above right=of zero] (b) {\(\mathbf b\)};
 \node[degree,above=28mm of zero] (jump) {\(\mathbf0'\)};
 \draw[order] (zero) -- (a);
 \draw[order] (zero) -- (b);
 \draw[order] (a) -- (jump);
 \draw[order] (b) -- (jump);
 \draw[jumpop,bend left=38] (zero) to
   node[right,font=\scriptsize] {jump \(\mathbf0\mapsto\mathbf0'\)} (jump);
 \node[font=\scriptsize] at ($(a)!0.5!(b)$) {incomparable};
\end{tikzpicture}
\caption{A schematic fragment of the degree order, not to scale
(apparatus A-126).  Here \(\mathbf a,\mathbf b\) are the c.e. degrees supplied
by Friedberg--Muchnik; the Kleene--Post theorem separately supplies an
unrestricted incomparable pair.  Solid arrows record only the displayed
strict comparisons; the dashed arrow is the established jump instance
\(\mathbf0\mapsto\mathbf0'\).  The figure asserts neither covers nor
distances, exhaustion of an interval, or the existence of a meet.}
\label{fig:p04-turing-degrees}
\end{figure}

%% Migration accounting for v1.1 lines 8898--9079:
%% - oracle computation, jump, relativized adequacy, the jump theorem, Post's
%%   theorem, and the convention bridge are all retained at source status;
%% - no source exercise occurs in this span;
%% - degree/priority/c.e.-structure entries are new registered content, kept
%%   distinct from the migrated source claims.
